\documentclass[11pt,a4paper]{article}

% Compile avec pdfLaTeX, XeLaTeX ou LuaLaTeX.
\usepackage{iftex}
\ifPDFTeX
  \usepackage[utf8]{inputenc}
  \usepackage[T1]{fontenc}
  \usepackage{lmodern}
\else
  \usepackage{fontspec}
\fi
\usepackage[french]{babel}
\usepackage[margin=2.5cm]{geometry}
\usepackage{array}
\usepackage{booktabs}
\usepackage{tabularx}
\usepackage{longtable}
\usepackage{float}
\usepackage{caption}
\usepackage{listings}
\usepackage{tikz}
\usetikzlibrary{arrows.meta,shapes.multipart}
\usepackage[hidelinks]{hyperref}

\setlength{\emergencystretch}{2em}
\newcolumntype{L}[1]{>{\raggedright\arraybackslash}p{#1}}
\newcolumntype{Y}{>{\raggedright\arraybackslash}X}
\newcommand{\code}[1]{\texttt{#1}}

\lstset{
  language=Java,
  basicstyle=\ttfamily\small,
  keywordstyle=\bfseries,
  commentstyle=\itshape,
  showstringspaces=false,
  columns=fullflexible,
  keepspaces=true,
  frame=single,
  framesep=4pt,
  xleftmargin=6pt,
  xrightmargin=6pt,
  breaklines=true,
  aboveskip=10pt,
  belowskip=8pt
}

\begin{document}

% ------------------------------------------------------------------ page de titre
\begin{titlepage}
\centering
{\large Programmation Java\par}
\vspace{0.3cm}
{\large Année 2026/2027\par}
\vfill
\rule{\textwidth}{0.6pt}\par
\vspace{0.7cm}
{\Huge\bfseries Gestion de stock\par}
\vspace{0.3cm}
{\Huge\bfseries d'une boutique solidaire\par}
\vspace{0.7cm}
{\Large Rapport final\par}
\vspace{0.7cm}
\rule{\textwidth}{0.6pt}\par
\vfill
% A remplacer par les noms des membres du groupe.
{\large
\begin{tabular}{c}
Prénom NOM \\
Prénom NOM \\
Prénom NOM \\
Prénom NOM \\
Prénom NOM \\
\end{tabular}\par}
\vfill
{\large Novembre 2026\par}
\end{titlepage}

\tableofcontents
\newpage

% ------------------------------------------------------------------ introduction
\section{Introduction}

Ce rapport présente le travail réalisé pendant les TP3 et TP4 du cours de Conception et
Programmation Objet. L'objectif était d'écrire en Java un programme en mode console qui gère le
stock d'une boutique solidaire. Cette boutique récupère des dons de particuliers et des invendus de
la grande distribution, puis les revend à prix modéré. Elle propose des articles culturels
d'occasion : des livres, des manuels scolaires et des magazines.

Le rapport suit l'ordre de notre travail. La section~\ref{sec:cahier} résume ce que demande
l'énoncé. La section~\ref{sec:conception} présente les classes et leurs relations. La
section~\ref{sec:realisation} explique les parties du code qui nous ont demandé le plus de
réflexion. La section~\ref{sec:verification} vérifie, question par question, que tout ce qui était
demandé est réalisé, puis décrit nos tests. La section~\ref{sec:limites} donne les limites du
programme.

% ------------------------------------------------------------------ cahier des charges
\section{Cahier des charges}
\label{sec:cahier}

\subsection{Les articles}

Tous les articles ont une description, un prix initial de vente et un nombre d'exemplaires. Chaque
type d'article ajoute ses propres informations (tableau~\ref{tab:articles}). Un manuel scolaire est
un livre particulier : il a tout ce qu'a un livre, avec une matière et un niveau scolaire en plus.

\begin{table}[H]
\centering
\begin{tabular}{@{}lll@{}}
\toprule
Type d'article & Numéro & Informations propres \\
\midrule
Livre & ISBN (13 chiffres) & nombre de pages \\
Manuel scolaire & ISBN (13 chiffres) & nombre de pages, matière, niveau scolaire \\
Magazine & ISSN (8 chiffres) & périodicité, date de publication \\
\bottomrule
\end{tabular}
\caption{Les trois types d'articles}
\label{tab:articles}
\end{table}

\subsection{Les prix}

Le prix d'un article dépend de la date (tableau~\ref{tab:prix}). Les livres et les manuels sont à
moitié prix pendant le mois d'avril. Pour un magazine, la réduction dépend de son ancienneté,
comptée dans l'unité de sa périodicité.

\begin{table}[H]
\centering
\begin{tabular}{@{}lll@{}}
\toprule
Article & Réduction de 50~\% & Réduction de 75~\% \\
\midrule
Livre, manuel & pendant le mois d'avril & jamais \\
Magazine hebdomadaire & publié depuis plus de 2 semaines & depuis plus de 4 semaines \\
Magazine mensuel & publié depuis plus de 2 mois & depuis plus de 4 mois \\
Magazine trimestriel & publié depuis plus de 2 trimestres & depuis plus de 4 trimestres \\
\bottomrule
\end{tabular}
\caption{Les réductions de prix}
\label{tab:prix}
\end{table}

\subsection{Les bons de dépôt}

À chaque dépôt d'articles, un bon de dépôt est établi. Il contient le numéro de téléphone du
déposant, la date du dépôt et, pour chaque article déposé, son numéro ISBN ou ISSN et le nombre
d'exemplaires. Le bon doit être numéroté et daté automatiquement. Selon le cas, l'article déposé
est ajouté au catalogue ou son nombre d'exemplaires est mis à jour.

\subsection{Les fichiers et le menu}

Les articles et les bons de dépôt sont sauvegardés dans deux fichiers texte dont le format est
imposé. Le programme principal charge ces fichiers au démarrage, propose un menu qui donne accès
aux fonctionnalités, puis sauvegarde les données avant de quitter.

\subsection{Les contraintes}

L'énoncé impose les noms des classes et des méthodes. Les données sont stockées dans des tableaux
de taille limitée. Les dates utilisent la classe \code{LocalDate}. Les fichiers sont écrits avec
\code{FileWriter} et lus avec \code{FileReader} et \code{BufferedReader}. Il faut aussi éviter de
mélanger les saisies, les affichages et les traitements dans les méthodes.

% ------------------------------------------------------------------ conception
\section{Conception}
\label{sec:conception}

\subsection{Vue d'ensemble}

Le programme compte huit classes, plus une classe \code{Tests}. La figure~\ref{fig:classes} montre
leurs relations et le tableau~\ref{tab:classes} donne le rôle de chacune.

\begin{figure}[H]
\centering
\begin{tikzpicture}[
  classe/.style={draw, rectangle split, rectangle split parts=2, align=center,
                 font=\small, text width=3.9cm, inner sep=4pt},
  herite/.style={-{Triangle[open, length=2.6mm, width=2.6mm]}},
  contient/.style={{Diamond[length=3.2mm, width=1.9mm]}-}]
\node[classe, text width=5.4cm] (article) at (0,0)
  {\textbf{Article} \textit{(abstraite)}\nodepart{second} description, prix, nbExemplaires};
\node[classe] (livre) at (-2.4,-2.5) {\textbf{Livre}\nodepart{second} isbn, nbPages};
\node[classe] (magazine) at (2.4,-2.5)
  {\textbf{Magazine}\nodepart{second} issn, periodicite,\\ datePublication};
\node[classe] (manuel) at (-2.4,-5.0) {\textbf{Manuel}\nodepart{second} matiere, niveauScolaire};
\node[classe] (etab) at (7.7,0) {\textbf{Etablissement}\nodepart{second} nom,\\ articles, bonsDepot};
\node[classe] (bon) at (7.7,-2.5)
  {\textbf{BonDepot}\nodepart{second} numero, telephone,\\ dateEmission, lignes};
\node[classe] (ligne) at (7.7,-5.0) {\textbf{LigneDepot}\nodepart{second} numero, nbExemplaires};
\draw[herite] (livre.north) -- (livre.north |- article.south);
\draw[herite] (magazine.north) -- (magazine.north |- article.south);
\draw[herite] (manuel.north) -- (livre.south);
\draw[contient] (etab.west) -- node[above, font=\scriptsize] {100 au plus} (article.east);
\draw[contient] (etab.south) -- node[right, font=\scriptsize] {100 au plus} (bon.north);
\draw[contient] (bon.south) -- node[right, font=\scriptsize] {20 au plus} (ligne.north);
\end{tikzpicture}
\caption{Les classes du programme. La flèche se lit « hérite de », le losange « contient un
tableau de ».}
\label{fig:classes}
\end{figure}

\begin{table}[H]
\centering
\begin{tabularx}{\textwidth}{@{}lY@{}}
\toprule
Classe & Rôle \\
\midrule
\code{Article} & Classe abstraite : ce qui est commun à tous les articles. \\
\code{Livre} & Article avec un ISBN et un nombre de pages. \\
\code{Manuel} & Livre avec une matière et un niveau scolaire. \\
\code{Magazine} & Article avec un ISSN, une périodicité et une date de publication. \\
\code{LigneDepot} & Une ligne d'un bon : un numéro d'article et un nombre d'exemplaires. \\
\code{BonDepot} & Un bon de dépôt : numéro, téléphone, date et lignes. \\
\code{Etablissement} & La boutique : le tableau des articles et le tableau des bons. \\
\code{Main} & Le programme principal : menu, saisies et messages. \\
\code{Tests} & Vérifications des méthodes sur des instances créées à la main. \\
\bottomrule
\end{tabularx}
\caption{Rôle des classes}
\label{tab:classes}
\end{table}

\subsection{La hiérarchie des articles}

La classe \code{Article} est abstraite : on ne crée jamais un article « tout court », on crée
toujours un livre, un manuel ou un magazine. Elle contient les trois attributs communs et les
méthodes qui ne dépendent pas du type : \code{placerApres}, \code{ajouter} et \code{retirer}.

Deux méthodes sont abstraites, parce que chaque type d'article les réalise à sa façon :
\code{getNumero}, qui renvoie l'ISBN ou l'ISSN, et \code{calculerPrix}. Grâce à cela, la classe
\code{Etablissement} manipule un seul tableau d'articles sans se demander de quel type ils sont.
Par exemple, \code{placerApres} est écrite une seule fois dans \code{Article} :

\begin{lstlisting}
public boolean placerApres(Article a) {
    return getNumero().compareTo(a.getNumero()) > 0;
}
\end{lstlisting}

\code{Manuel} hérite de \code{Livre}. Il n'a donc pas besoin de redéfinir \code{getNumero} ni
\code{calculerPrix} : la réduction du mois d'avril écrite dans \code{Livre} s'applique aussi aux
manuels.

Pour les attributs, \code{description} et \code{nbExemplaires} sont privés. Seul \code{prix} est
\code{protected}, parce que les sous-classes en ont besoin pour calculer le prix du jour.

\subsection{Les bons de dépôt}

Une \code{LigneDepot} contient le numéro d'un article et un nombre d'exemplaires. Un
\code{BonDepot} contient un tableau de vingt lignes au plus et un compteur \code{nbLignes}.

La numérotation automatique repose sur un attribut \code{static}, donc partagé par tous les bons.
Le premier constructeur sert pour un nouveau bon : il prend le numéro suivant et la date du jour.

\begin{lstlisting}
private static int prochainNumero = 1;

public BonDepot(String telephone) {
    this.numero = prochainNumero;
    prochainNumero++;
    this.telephone = telephone;
    this.dateEmission = LocalDate.now();
    this.lignes = new LigneDepot[20];
    this.nbLignes = 0;
}
\end{lstlisting}

Le second constructeur sert à recréer un bon lu dans le fichier, avec son numéro et sa date
d'origine. Il avance le compteur si nécessaire, pour que le prochain bon créé porte toujours un
numéro plus grand que tous les bons déjà enregistrés :

\begin{lstlisting}
if (numero >= prochainNumero) {
    prochainNumero = numero + 1;
}
\end{lstlisting}

\subsection{La classe Etablissement}

La boutique contient deux tableaux de cent cases, chacun avec un compteur : \code{articles} avec
\code{nbArticles}, et \code{bonsDepot} avec \code{nbBonsDepot}. Le tableau des articles est trié
par numéro ISBN ou ISSN, le tableau des bons par date d'émission, le plus récent en dernier.

L'énoncé demande plusieurs méthodes qui portent le même nom \code{ajouter}. Elles se distinguent
par leurs paramètres (tableau~\ref{tab:ajouter}).

\begin{table}[H]
\centering
\begin{tabularx}{\textwidth}{@{}L{8.4cm}Y@{}}
\toprule
Paramètres de \code{ajouter} & Rôle \\
\midrule
description, prix, nbExemplaires, isbn, nbPages & ajoute un livre \\
les mêmes, puis matiere, niveauScolaire & ajoute un manuel \\
description, prix, nbExemplaires, issn, periodicite, datePublication & ajoute un magazine \\
numero, quantite & augmente le stock d'un article \\
bon & enregistre un bon de dépôt \\
\bottomrule
\end{tabularx}
\caption{Les cinq méthodes \code{ajouter} de la classe \code{Etablissement}}
\label{tab:ajouter}
\end{table}

Ces méthodes renvoient un booléen : vrai si l'opération a réussi, faux sinon (tableau plein, numéro
déjà utilisé, article inconnu). C'est le programme principal qui affiche le message correspondant.

% ------------------------------------------------------------------ realisation
\section{Réalisation}
\label{sec:realisation}

\subsection{Insertion dans un tableau trié}

Le tableau des articles doit rester trié par numéro. Plutôt que de trier tout le tableau après
chaque ajout, nous insérons le nouvel article directement à sa place.

\begin{lstlisting}
private boolean insererArticle(Article a) {
    if (nbArticles == articles.length || rechercher(a.getNumero()) != null) {
        return false;
    }
    // on decale vers la droite les articles qui doivent etre apres le nouveau
    int i = nbArticles;
    while (i > 0 && articles[i - 1].placerApres(a)) {
        articles[i] = articles[i - 1];
        i--;
    }
    articles[i] = a;
    nbArticles++;
    return true;
}
\end{lstlisting}

On commence par vérifier que le tableau n'est pas plein et que le numéro n'existe pas déjà. Ensuite
on part de la première case libre, à la fin. Tant que l'article qui précède doit être placé après
le nouveau, on le recopie une case plus loin et on recule. Quand la boucle s'arrête, la case
\code{i} est libre et c'est la bonne place.

Par exemple, si le tableau contient les numéros \code{00208245}, \code{9782210446694} et
\code{9782707356741} et que l'on ajoute \code{9782070612758}, les deux derniers articles sont
décalés d'une case et le nouveau prend la deuxième place.

Les bons de dépôt sont rangés de la même manière dans \code{insererBon}, en comparant les dates
avec \code{isAfter}. Deux bons de la même date restent dans leur ordre d'arrivée.

\subsection{Calcul des prix}

Pour un livre, le calcul tient en un test sur le mois de la date. Pour un magazine, nous calculons
deux dates limites à partir de la date de publication, puis nous comparons :

\begin{lstlisting}
public double calculerPrix(LocalDate date) {
    LocalDate limite50;
    LocalDate limite75;
    if (periodicite.equals("hebdomadaire")) {
        limite50 = datePublication.plusWeeks(2);
        limite75 = datePublication.plusWeeks(4);
    } else if (periodicite.equals("mensuel")) {
        limite50 = datePublication.plusMonths(2);
        limite75 = datePublication.plusMonths(4);
    } else if (periodicite.equals("trimestriel")) {
        // 1 trimestre = 3 mois
        limite50 = datePublication.plusMonths(6);
        limite75 = datePublication.plusMonths(12);
    } else {
        return prix;
    }
    if (date.isAfter(limite75)) {
        return prix / 4;
    }
    if (date.isAfter(limite50)) {
        return prix / 2;
    }
    return prix;
}
\end{lstlisting}

L'énoncé dit « publié depuis \emph{plus de} deux semaines ». Nous l'avons compris au sens strict,
d'où \code{isAfter} : le jour exact des deux semaines, le magazine est encore à plein tarif, et la
réduction commence le lendemain.

La méthode reçoit la date en paramètre au lieu d'utiliser directement la date du jour. Cela permet
de tester le calcul à n'importe quelle date. Pour l'affichage, \code{lister} lui passe
\code{LocalDate.now()}.

Exemple : le magazine \emph{Geo} de nos données est un mensuel publié le 1\ier{} août 2024. Il a
dépassé quatre mois depuis longtemps, donc son prix du jour est de $6{,}50 / 4 = 1{,}625$ euro.

\subsection{Consultation : les trois méthodes lister}

\begin{itemize}
\item \code{lister()} affiche les articles par nombre d'exemplaires croissant, avec leur prix du
jour. Le tri est un tri à bulles fait sur une copie du tableau, pour ne pas casser le tri par
numéro du tableau d'origine.
\item \code{lister(telephone)} affiche les bons d'un client. Comme le tableau des bons est déjà
trié par date, il suffit de le parcourir.
\item \code{lister(numero, debut, fin)} affiche les bons qui contiennent un article sur une
période. Les deux dates limites sont incluses, ce qui s'écrit
\code{!date.isBefore(debut) \&\& !date.isAfter(fin)}.
\end{itemize}

\subsection{Sauvegarde dans des fichiers texte}

\subsubsection*{Le format}

Dans le fichier des articles, chaque article occupe deux lignes : son numéro, puis ses informations
séparées par \code{" : "}. Dans le fichier des bons, chaque bon commence par son numéro, puis une
ligne d'en-tête (téléphone, date, nombre d'articles), puis une ligne par article déposé.

\begin{lstlisting}[language={}]
00208245
Geo : 6.5 : 20 : mensuel : 2024-08-01
9782210446694
Nouvelle grammaire du college : 38.0 : 31 : 512 : Francais : College
9782707356741
La maison vide : 25.0 : 10 : 752
\end{lstlisting}

\begin{lstlisting}[language={}]
1
061236547890 : 2025-10-10 : 2
10 : 9782210446694
20 : 9782707356741
\end{lstlisting}

\subsubsection*{L'écriture}

Chaque classe a une méthode \code{versFichier} qui renvoie son texte. Dans la hiérarchie des
articles, chaque classe écrit seulement ce qu'elle ajoute et appelle celle de sa classe mère avec
\code{super} :

\begin{lstlisting}
// dans Article
return getNumero() + System.lineSeparator()
        + description + " : " + prix + " : " + nbExemplaires;
// dans Livre
return super.versFichier() + " : " + nbPages;
// dans Manuel
return super.versFichier() + " : " + matiere + " : " + niveauScolaire;
// dans Magazine
return super.versFichier() + " : " + periodicite + " : " + datePublication;
\end{lstlisting}

Dans \code{Etablissement}, \code{versFichierArticles} et \code{versFichierDepots} parcourent leur
tableau et écrivent chaque élément avec un \code{FileWriter}. Le saut de ligne est écrit entre deux
éléments seulement, pas à la fin, pour que le fichier ne contienne aucune ligne vide.

\subsubsection*{La relecture}

Le fichier ne dit pas de quel type est un article. Nous le déduisons du nombre d'informations de la
deuxième ligne, obtenu avec \code{split} :

\begin{lstlisting}
String[] infos = fichier.readLine().split(" : ");
String description = infos[0];
double prix = Double.parseDouble(infos[1]);
int nbExemplaires = Integer.parseInt(infos[2]);
// le nombre d'infos donne le type : 4 = livre, 6 = manuel, 5 = magazine
if (infos.length == 4) {
    ajouter(description, prix, nbExemplaires, numero, Integer.parseInt(infos[3]));
} else if (infos.length == 6) {
    ajouter(description, prix, nbExemplaires, numero, Integer.parseInt(infos[3]),
            infos[4], infos[5]);
} else {
    ajouter(description, prix, nbExemplaires, numero, infos[3],
            LocalDate.parse(infos[4]));
}
\end{lstlisting}

Comme la relecture passe par les méthodes \code{ajouter}, les articles relus sont insérés à leur
place : le tableau est trié même si le fichier ne l'était pas.

Pour les bons, il y a un piège. Enregistrer un bon augmente le stock des articles déposés. Mais le
stock écrit dans le fichier des articles tient déjà compte de tous les dépôts passés. À la
relecture, il ne faut donc pas appeler \code{ajouter(bon)}, sinon chaque dépôt serait compté une
deuxième fois. \code{depuisFichierDepots} appelle seulement \code{insererBon}, qui range le bon
dans le tableau sans toucher au stock.

Enfin, comme \code{" : "} sert de séparateur, un texte qui contiendrait ce caractère casserait la
relecture. Le programme principal refuse donc le caractère « : » dans tous les textes saisis.

\subsection{Programme principal}

La classe \code{Main} crée la boutique, charge les deux fichiers, puis affiche le menu dans une
boucle jusqu'au choix 0 (tableau~\ref{tab:menu}). Avant de quitter, elle sauvegarde les deux
fichiers.

\begin{table}[H]
\centering
\begin{tabularx}{\textwidth}{@{}clY@{}}
\toprule
Choix & Action & Méthode de \code{Etablissement} \\
\midrule
1 & Afficher la boutique & \code{afficher} \\
2 & Lister les articles par stock & \code{lister()} \\
3 & Rechercher un article & \code{rechercher} \\
4 & Ajouter un article (livre, manuel ou magazine) & \code{ajouter}, trois versions \\
5 & Augmenter un stock & \code{ajouter(numero, quantite)} \\
6 & Diminuer un stock & \code{retirer} \\
7 & Enregistrer un bon de dépôt & \code{ajouter(bon)} \\
8 & Lister les bons d'un client & \code{lister(telephone)} \\
9 & Lister les bons d'un article sur une période & \code{lister(numero, debut, fin)} \\
0 & Quitter & \code{versFichierArticles}, \code{versFichierDepots} \\
\bottomrule
\end{tabularx}
\caption{Le menu}
\label{tab:menu}
\end{table}

Nous avons séparé les rôles. Toutes les saisies au clavier sont dans \code{Main}. Les autres
classes font les traitements et renvoient un résultat, et \code{Main} affiche le message.

Les saisies sont contrôlées par quatre petites méthodes : \code{lireTexte}, \code{lireEntier},
\code{lirePrix} et \code{lireDate}. Chacune redemande la valeur tant qu'elle n'est pas correcte
(tableau~\ref{tab:saisies}). Une cinquième, \code{lireNumero}, enlève les espaces : un ISBN tapé
\code{978 2 7073 5674 1} ou \code{9782707356741} désigne ainsi le même article.

\begin{table}[H]
\centering
\begin{tabularx}{\textwidth}{@{}L{6.2cm}Y@{}}
\toprule
Saisie & Réaction du programme \\
\midrule
Texte vide & redemandé \\
Texte contenant « : » & redemandé, c'est le séparateur des fichiers \\
Lettres à la place d'un nombre & redemandé \\
Nombre hors limites (choix du menu, quantité, prix) & redemandé, avec les limites \\
Date mal écrite & redemandée, au format AAAA-MM-JJ \\
Périodicité inconnue & redemandée \\
Numéro d'article déjà utilisé & ajout refusé \\
Retrait supérieur au stock & refusé, le stock ne change pas \\
\bottomrule
\end{tabularx}
\caption{Contrôle des saisies}
\label{tab:saisies}
\end{table}

Pour le bon de dépôt (choix 7), le programme demande le téléphone du déposant, puis chaque article
et sa quantité. Si l'article n'est pas encore au catalogue, il est créé à ce moment-là avec zéro
exemplaire, et c'est le bon qui apporte les exemplaires déposés. On retrouve ainsi les deux cas de
l'énoncé : soit l'article est ajouté, soit son nombre d'exemplaires est mis à jour.

% ------------------------------------------------------------------ verification
\section{Vérification}
\label{sec:verification}

\subsection{Conformité à l'énoncé}

Le tableau~\ref{tab:conformite} reprend l'énoncé question par question. Toutes les questions du TP3
et du TP4 sont réalisées.

\begin{longtable}{@{}L{1.6cm}L{5.3cm}L{6.4cm}L{1.0cm}@{}}
\caption{Conformité à l'énoncé}\label{tab:conformite}\\
\toprule
Question & Ce qui est demandé & Ce qui est réalisé & État \\
\midrule
\endfirsthead
\toprule
Question & Ce qui est demandé & Ce qui est réalisé & État \\
\midrule
\endhead
\bottomrule
\endfoot
\multicolumn{4}{@{}l}{\textbf{TP3}} \\
1 (a) & Classes \code{Article}, \code{Livre}, \code{Manuel}, \code{Magazine} : attributs,
constructeurs, \code{toString} & Les quatre classes, avec héritage et \code{toString} dans chacune
& Fait \\
1 (b) & Classe \code{LigneDepot} & Numéro, nombre d'exemplaires, constructeur, \code{toString}
& Fait \\
2 & Classe \code{BonDepot} : téléphone, date, lignes dans un tableau, numéro unique, affichage
& Tableau de 20 lignes, compteur \code{static}, méthode \code{afficher} & Fait \\
3 & Classe \code{Etablissement} : nom, tableaux des articles et des bons, affichage
& Deux tableaux de 100 cases avec compteurs, méthode \code{afficher} & Fait \\
4 (a) & \code{getNumero}, \code{placerApres}, \code{ajouter}, \code{retirer}, \code{calculerPrix}
& Dans \code{Article}, les deux méthodes abstraites sont définies dans \code{Livre} et
\code{Magazine} & Fait \\
4 (b) & \code{ajouter} pour un livre, un manuel, un magazine & Trois méthodes \code{ajouter}, avec
insertion triée par numéro & Fait \\
4 (c) & \code{rechercher}, \code{ajouter} et \code{retirer} avec un numéro et une quantité
& \code{rechercher} renvoie l'article ou \code{null}, les deux autres renvoient un booléen
& Fait \\
5 (a) & \code{ajouterLigne} dans \code{BonDepot} & Ajoute une ligne, refusée si le bon est plein ou
la quantité nulle & Fait \\
5 (b) & \code{ajouter} un bon dans \code{Etablissement}, bons numérotés et datés automatiquement
& Insertion triée par date, mise à jour du stock, numéro et date fixés par le constructeur
& Fait \\
\addlinespace
\multicolumn{4}{@{}l}{\textbf{TP4}} \\
1 & Trois méthodes \code{lister} & Articles par stock croissant avec prix du jour, bons d'un
client, bons d'un article sur une période & Fait \\
2 (a) & \code{versFichier} dans les articles, \code{LigneDepot} et \code{BonDepot}
& Une méthode par classe, enchaînées avec \code{super} pour les articles & Fait \\
2 (b) & \code{versFichierArticles}, \code{depuisFichierArticles} & Écriture avec \code{FileWriter},
lecture avec \code{FileReader} et \code{BufferedReader} & Fait \\
2 (c) & \code{versFichierDepots}, \code{depuisFichierDepots} & Mêmes classes, la relecture ne
modifie pas le stock & Fait \\
3 & Programme principal : instance, chargement, menu, sauvegarde & Classe \code{Main}, menu de neuf
choix & Fait \\
\addlinespace
\multicolumn{4}{@{}l}{\textbf{Consignes générales}} \\
& Noms de classes et de méthodes imposés & Tous respectés & Fait \\
& Stockage dans des tableaux de taille limitée & 100 articles, 100 bons, 20 lignes par bon
& Fait \\
& Articles triés par numéro, bons triés par date & Insertion à la bonne place dans les deux
tableaux & Fait \\
& Dates avec \code{LocalDate} & Date de publication, date d'émission, calcul des prix & Fait \\
& Format des deux fichiers & Fichiers produits identiques à l'exemple de l'énoncé & Fait \\
& Dépôt : article ajouté ou stock mis à jour & Choix 7 : un article inconnu est créé, puis le bon
met le stock à jour & Fait \\
& Éviter de mélanger saisies, affichages et traitements & Saisies dans \code{Main} seulement, les
traitements renvoient un booléen & Fait \\
& Tester les méthodes en créant des instances & Classe \code{Tests}, 21 vérifications & Fait \\
\end{longtable}

Une remarque sur la séparation des affichages et des traitements : les méthodes \code{afficher} et
\code{lister} affichent, puisque c'est leur rôle. Les quatre méthodes de fichiers affichent aussi
un message quand un fichier est absent ou ne peut pas être écrit. Nous avons gardé ce message à cet
endroit pour que le programme principal reste court.

\subsection{Choix d'interprétation}

Sur quelques points, l'énoncé laisse un choix. Voici ceux que nous avons faits.

\begin{itemize}
\item « Plus de deux semaines » est compris au sens strict : la réduction commence le lendemain du
jour des deux semaines.
\item Un trimestre vaut trois mois, comptés à partir de la date de publication.
\item Pour la période de \code{lister}, les deux dates limites sont incluses.
\item Le prix du jour n'est pas arrondi : il est affiché tel qu'il est calculé.
\item Les fichiers sont écrits sans ligne vide et sans saut de ligne final, comme dans l'exemple de
l'énoncé.
\item À la relecture, les bons ne modifient pas le stock.
\end{itemize}

\subsection{Tests}

\subsubsection*{La classe Tests}

La classe \code{Tests} crée des instances à la main et compare le résultat de chaque méthode avec
la valeur attendue. Elle affiche \code{OK} ou \code{ERREUR} pour chaque vérification, puis le
nombre d'erreurs. Elle contient 21 vérifications (tableau~\ref{tab:tests}), toutes réussies.

\begin{table}[H]
\centering
\begin{tabularx}{\textwidth}{@{}L{3.6cm}Yc@{}}
\toprule
Sujet & Vérifications & Nombre \\
\midrule
Prix des livres et des manuels & Plein tarif en janvier, moitié prix en avril & 3 \\
Prix des magazines & Pour chaque périodicité : le jour du seuil, le lendemain, et le lendemain
du second seuil & 9 \\
Stock d'un article & Retrait possible, retrait trop grand refusé, ajout de zéro refusé & 3 \\
Catalogue & Numéro déjà utilisé refusé, recherche d'un article connu et inconnu, stock d'un
article inconnu refusé & 4 \\
Bon de dépôt & Ligne de zéro exemplaire refusée, stock augmenté après le dépôt & 2 \\
\bottomrule
\end{tabularx}
\caption{Les vérifications de la classe \code{Tests}}
\label{tab:tests}
\end{table}

\subsubsection*{Les scénarios au clavier}

Nous avons ensuite utilisé le programme comme le ferait le gérant de la boutique
(tableau~\ref{tab:scenarios}).

\begin{table}[H]
\centering
\begin{tabularx}{\textwidth}{@{}L{7.0cm}Y@{}}
\toprule
Scénario & Résultat \\
\midrule
Premier lancement sans aucun fichier & Boutique vide, fichiers créés à la sortie \\
Chargement des fichiers d'exemple, puis sortie & Fichiers réécrits à l'identique \\
Ajout d'un livre, d'un manuel et d'un magazine dans le désordre & Catalogue trié par numéro \\
Augmentation puis diminution d'un stock & Stock mis à jour \\
Retrait supérieur au stock, article inconnu & Refusé, stock inchangé \\
Bon de dépôt sur deux articles du catalogue & Bon numéro 2, date du jour, stocks augmentés \\
Bon de dépôt avec un article inconnu & Article créé, stock égal à la quantité déposée \\
Fermeture du programme puis relance & Mêmes données, le bon suivant porte le numéro 3 \\
Saisies incorrectes (tableau~\ref{tab:saisies}) & Message et nouvelle demande, pas d'arrêt du
programme \\
Numéro tapé avec des espaces & Le même article est retrouvé \\
\bottomrule
\end{tabularx}
\caption{Scénarios testés avec le menu}
\label{tab:scenarios}
\end{table}

Enfin, le programme compile sans aucun avertissement, même avec l'option \code{-Xlint:all} du
compilateur.

% ------------------------------------------------------------------ limites
\section{Limites et améliorations possibles}
\label{sec:limites}

Le programme a des limites que nous connaissons.

\begin{itemize}
\item Le catalogue est limité à 100 articles, la boutique à 100 bons et un bon à 20 lignes. Ces
tailles sont fixées à la création des tableaux.
\item On ne peut ni supprimer un article, ni modifier sa description ou son prix.
\item Le caractère « : » est interdit dans les textes, parce qu'il sert de séparateur dans les
fichiers.
\item La longueur des numéros n'est pas vérifiée : le programme accepte un ISBN qui n'a pas 13
chiffres.
\item Les données sont sauvegardées seulement quand on quitte par le choix 0.
\item Les fichiers sont toujours écrits par le programme. Un fichier modifié à la main et mal formé
arrête le programme au démarrage.
\item Les données d'exemple sont écrites sans accents, parce que l'affichage des accents dépend de
la console utilisée.
\end{itemize}

Avec plus de temps, nous ajouterions la suppression et la modification d'un article, la vérification
de la longueur des numéros ISBN et ISSN, et une interface graphique à la place du menu.

% ------------------------------------------------------------------ conclusion
\section{Conclusion}

Le programme répond à toutes les questions du TP3 et du TP4 : il gère les trois types d'articles,
calcule les prix avec les réductions, enregistre les bons de dépôt et conserve les données d'une
exécution à l'autre dans deux fichiers texte au format demandé.

Nous avons cherché à garder un code simple, qui utilise seulement des tableaux et les classes vues
en cours. Ce projet nous a surtout permis de pratiquer l'héritage et les classes abstraites, de
voir l'intérêt d'une méthode redéfinie dans chaque sous-classe, et d'apprendre à lire et à écrire
dans un fichier texte.

% ------------------------------------------------------------------ annexe
\appendix
\section{Utilisation du programme}

Le projet est un projet Maven, écrit pour Java 26. Les sources Java sont rangées dans
\code{src/main/java}.

\begin{itemize}
\item Dans NetBeans, on ouvre le projet puis on lance \emph{Run Project}. Le menu apparaît dans la
fenêtre \emph{Output}.
\item Les fichiers \code{articles.txt} et \code{depots.txt} sont lus et écrits à la racine du
projet. S'ils n'existent pas, la boutique démarre vide et ils sont créés à la sortie.
\item Pour lancer les vérifications, on fait un clic droit sur \code{Tests.java}, puis
\emph{Run File}.
\end{itemize}

\end{document}
